\documentclass[11pt]{letter}

\usepackage[margin=1in]{geometry}
\usepackage{amsmath}
\usepackage{amssymb}
\usepackage{hyperref}
\hypersetup{colorlinks=true,linkcolor=blue,citecolor=blue,urlcolor=blue}

\signature{Conrard Giresse Tetsassi Feugmo}
\address{Department of Chemistry and \\ Department of Physics \& Astronomy \\
University of Waterloo \\ 200 University Avenue West \\ Waterloo, ON N2L 3G1, Canada \\
\texttt{cgtetsas@uwaterloo.ca}}

\begin{document}
\setlength{\parskip}{0.4ex}

\begin{letter}{The Editors \\ \textit{npj Materials Degradation} \\ Nature Portfolio}

\opening{Dear Editors,}

We are pleased to submit our manuscript, \textbf{``Sparse data limit
what a corrosion model can predict,''} for consideration as an Article
in \textit{npj Materials Degradation}.

Mechanistic corrosion models are routinely fitted with four to six
parameters to a handful of measurements and then extrapolated across
service lifetimes. Whether those parameters are determined by the data is
almost never tested, and our central result is that, for a
representative case, several of them are not, with consequences large enough to change a lifetime assessment. We establish this for the duplex
oxide grown on Nb-stabilized AISI~347 (X6CrNiNb18-10) in simulated
boiling-water-reactor hydrothermal water (240\,$^\circ$C, 7\,MPa,
0.4\,ppm dissolved O$_2$), using a differentiable, physics-informed
inversion of the point defect model (PDM). It is well suited to \textit{npj Materials Degradation}:

\begin{itemize}
\item \textbf{The first identifiability characterization of a PDM fit.}
Profile likelihood, a 1000-member bootstrap, a prior-leverage test and a
residual-leverage decomposition are built into the inversion rather than
added afterward. They show that of the five parameters of the accepted
model, three exposure times leave one undetermined and pin a second only
to within an order of magnitude, and that 95\% of the fit statistic is
carried by three chromium points, and 80\% by a single measurement. These tools are mature in
adjacent fields but, to our knowledge, have not been brought to the PDM;
the diagnosis transfers to any sparsely sampled degradation model.

\item \textbf{Consequences large enough to matter, stated at the
resolution the data support.} The identified kinetics fit the
calibration data as well as an empirical power law but extrapolate
differently: at ten years the power law lies a factor of 3.5 to 6.3
above the mechanistic barrier layer and outside its uncertainty band
under every error treatment we apply, including one that widens that
band from 447--704\,nm to roughly 340--1040\,nm. We are deliberate that
the \emph{separation} is the result and the point estimate it brackets
(535\,nm) is not: an analysis whose thesis is that sparse data
under-determine parameters should not then quote a lifetime number as
though it were determined. A designed exposure beyond roughly
3000\,hours would resolve the separation, a concrete recommendation for lifetime assessment and the reason the re-analysis is worth doing
on data already in the literature.

\item \textbf{A diagnosed and remedied failure mode.} We document a
physics-informed inversion failure specific to sparse experimental data, the data loss satisfied while the governing equation is not, and remedy it with a forward-solve self-consistency check, a
cautionary and constructive result for the growing use of
physics-informed learning in this field.

\item \textbf{A negative result on the PDM's own constant-field
closure.} Solving Poisson's equation self-consistently with the defect
transport, instead of imposing the field, shows the closure is not
supported at the defect concentrations the model implies: the implied Debye length is about 0.06\,nm, several times smaller than an interatomic spacing. This does not disturb the fitted kinetics, which
identify $b_3$ rather than the field itself, but it bounds how the
derived field strengths may be read, and we give it its own section
rather than a footnote.

\item \textbf{A first mechanistic description for this alloy, with its limits stated.} Prior characterization of AISI~347 oxide in BWR water stopped
at empirical per-element power-law fits with no shared physical parameter
set; we connect the Cr-rich barrier and Fe-rich outer layers through a
single kinetic description. We are equally explicit that the model
structure is not itself Nb-specific and that the neural inversion is
validated against an exact classical reference rather than required by the problem.
\end{itemize}

The manuscript is original, has not been published elsewhere, and is not
under consideration by any other journal. The author declares no
competing interests. The work involves no human or animal subjects. All
code, the digitized data, and a committed artefact behind every number
in the manuscript are available in a public repository, together with a
generated map from each figure and table to the command that produces
it, so that any reported value can be re-derived without re-running the
full pipeline.
Thank you for considering our submission.

\closing{Sincerely,}

\end{letter}

\end{document}
